\chapter{Sequences}\label{ch:g12:seq}

A \omterm{def:g12:seq:sequence}{sequence} is a list of \omterm{def:g10:numbers:sets}{real numbers} indexed by the \omterm{def:g10:numbers:sets}{natural numbers}. \omterm{def:g12:seq:sequence}{Sequences}
model discrete evolutions --- populations counted year by year, balances of a
bank account, successive \omterm{def:g10:numbers:approx}{approximations} of a number --- and their limits are
the first serious encounter with the infinite. This chapter sets up the
vocabulary, the principle of induction, and the fundamental convergence
theorems.

\section{Reasoning by induction}

\begin{theorem}[Principle of induction]\label{thm:g12:seq:induction}
Let $P(n)$ be a statement depending on an \omterm{def:g10:numbers:sets}{integer} $n$, and let $n_0 \in \N$.
If
\begin{enumerate}
  \item \emph{(base case)} $P(n_0)$ is true, and
  \item \emph{(inductive step)} for every $n \geq n_0$, $P(n)$ implies $P(n+1)$,
\end{enumerate}
then $P(n)$ is true for every $n \geq n_0$.
\end{theorem}

\begin{proof}
Suppose, for contradiction, that the set $A$ of \omterm{def:g10:numbers:sets}{integers} $n \geq n_0$ for
which $P(n)$ is false is nonempty. Then $A$ has a smallest element
$m$.\footnote{Every nonempty subset of $\N$ has a smallest element; this
property of $\N$ is taken as an axiom.} Since $P(n_0)$ is true, $m > n_0$, so
$m - 1 \geq n_0$ and $m-1 \notin A$, i.e.\ $P(m-1)$ is true. The inductive
step applied to $n = m-1$ then shows that $P(m)$ is true, contradicting
$m \in A$.
\end{proof}

\begin{example}\label{ex:g12:seq:bernoulli}
Let us prove \emph{Bernoulli's inequality}: for every real $a > 0$ and every
$n \in \N$,
\[
(1+a)^n \geq 1 + na.
\]
\emph{Base case.} For $n = 0$, both sides equal $1$.
\emph{Inductive step.} Suppose $(1+a)^n \geq 1+na$ for some $n \in \N$. Since
$1 + a > 0$, multiplying both sides by $1+a$ preserves the inequality:
\[
(1+a)^{n+1} \geq (1+na)(1+a) = 1 + (n+1)a + na^2 \geq 1 + (n+1)a .
\]
By induction, the inequality holds for all $n \in \N$.
\end{example}

\begin{method}[Writing an induction proof]\label{met:g12:seq:induction}
Always make the statement $P(n)$ explicit before starting. A complete proof
has three visible parts: the base case, the inductive step (``assume $P(n)$;
we prove $P(n+1)$''), and the conclusion invoking the principle of induction.
The most common error is to prove the inductive step without ever using the
hypothesis $P(n)$: if that happens, either the proof is wrong or induction was
not needed.
\end{method}

\section{Vocabulary of sequences}

\begin{definition}[Sequence]\label{def:g12:seq:sequence}
A \emph{sequence}\index{sequence} is a \omterm{def:g11:func:function}{function} $u \colon \N \to \R$ (or from
$\{n \in \N : n \geq n_0\}$ to $\R$). The \omterm{def:g10:functions:function}{image} of $n$ is written $u_n$, and
the sequence itself $(u_n)_{n\in\N}$ or simply $(u_n)$.
\end{definition}

A \omterm{def:g12:seq:sequence}{sequence} may be defined \emph{explicitly}, by a formula $u_n = f(n)$, or by
\emph{recurrence}, by its first term and a relation $u_{n+1} = f(u_n)$.

\begin{definition}[Monotonicity]\label{def:g12:seq:monotonic}
A \omterm{def:g12:seq:sequence}{sequence} $(u_n)$ is \emph{increasing}\index{sequence!increasing} if
$u_{n+1} \geq u_n$ for all $n$, \emph{\omterm{def:g10:functions:variations}{decreasing}} if $u_{n+1} \leq u_n$ for
all $n$, and \emph{monotonic} if it is increasing or \omterm{def:g10:functions:variations}{decreasing}. It is
\emph{strictly} increasing (resp.\ \omterm{def:g10:functions:variations}{decreasing}) when the inequalities are
strict.
\end{definition}

\begin{method}[Studying the monotonicity of a sequence]\label{met:g12:seq:monotonicity}
Three standard techniques:
\begin{enumerate}
  \item study the sign of $u_{n+1} - u_n$;
  \item if all terms are positive, compare $\dfrac{u_{n+1}}{u_n}$ to $1$;
  \item if $u_n = f(n)$ with $f$ defined on $\intco{0}{+\infty}$, use the
        variations of $f$.
\end{enumerate}
\end{method}

\begin{definition}[Bounded sequence]\label{def:g12:seq:bounded}
A \omterm{def:g12:seq:sequence}{sequence} $(u_n)$ is \emph{bounded above}\index{sequence!bounded} if there
exists $M \in \R$ with $u_n \leq M$ for all $n$; \emph{bounded below} if there
exists $m \in \R$ with $u_n \geq m$ for all $n$; and \emph{bounded} if both
hold.
\end{definition}

\subsection{Arithmetic and geometric sequences}

\begin{definition}[Arithmetic and geometric sequences]\label{def:g12:seq:arith-geom}
A \omterm{def:g12:seq:sequence}{sequence} $(u_n)$ is \emph{arithmetic}\index{sequence!arithmetic} with \omterm{def:g11:seq:arithmetic}{common
difference} $r$ if $u_{n+1} = u_n + r$ for all $n$, and
\emph{geometric}\index{sequence!geometric} with \omterm{def:g11:seq:geometric}{common ratio} $q$ if
$u_{n+1} = q\,u_n$ for all $n$.
\end{definition}

\begin{proposition}[Explicit form and sums]\label{prop:g12:seq:arith-geom}
Let $n \in \N$.
\begin{enumerate}
  \item If $(u_n)$ is \omterm{def:g12:seq:arith-geom}{arithmetic} with \omterm{def:g11:seq:arithmetic}{common difference} $r$, then
        $u_n = u_0 + nr$ and
        \[
        u_0 + u_1 + \dots + u_n = (n+1)\,\frac{u_0 + u_n}{2}.
        \]
  \item If $(u_n)$ is \omterm{def:g12:seq:arith-geom}{geometric} with \omterm{def:g11:seq:geometric}{common ratio} $q \neq 1$, then
        $u_n = u_0\, q^n$ and
        \[
        u_0 + u_1 + \dots + u_n = u_0\,\frac{1 - q^{n+1}}{1 - q}.
        \]
\end{enumerate}
\end{proposition}

\begin{proof}
The explicit forms follow by immediate inductions. For the \omterm{def:g12:seq:arith-geom}{arithmetic} sum,
write $S = u_0 + \dots + u_n$ and add the same sum written in reverse order:
each of the $n+1$ column sums equals $u_0 + u_n$, so $2S = (n+1)(u_0+u_n)$.
For the \omterm{def:g12:seq:arith-geom}{geometric} sum, compute $S - qS$: all terms cancel in pairs except the
first and the last, so $(1-q)S = u_0(1 - q^{n+1})$.
\end{proof}

\section{Limit of a sequence}

\begin{definition}[Convergent sequence]\label{def:g12:seq:limit}
A \omterm{def:g12:seq:sequence}{sequence} $(u_n)$ \emph{converges}\index{limit!of a sequence} to the \omterm{def:g10:numbers:sets}{real
number} $\ell$ if every open \omterm{def:g10:numbers:interval}{interval} containing $\ell$ contains all the terms
$u_n$ from some index on. We then write $\lim\limits_{n\to+\infty} u_n = \ell$.

Equivalently: for every $\varepsilon > 0$, there exists $N \in \N$ such that
for all $n \geq N$, $\abs{u_n - \ell} \leq \varepsilon$.
\end{definition}

\begin{omfigure}
\begin{tikzpicture}
\begin{axis}[omaxis,
  width=11cm, height=5.6cm,
  xlabel={$n$}, ylabel={$u_n$},
  xmin=0, xmax=16.5, ymin=0.7, ymax=2.9,
  xtick={4}, xticklabels={$N$}, ytick={2}, yticklabels={$\ell$}]
  \fill[omDef!20] (axis cs:0,1.75) rectangle (axis cs:16.5,2.25);
  \addplot[omProp, dashed, domain=0:16.5] {2};
  \addplot[omDef, only marks, mark=*, mark size=1.5pt,
    samples at={1,...,16}] {2 + (-1)^x/x};
  \draw[black, dashed, thin] (axis cs:4,0.7) -- (axis cs:4,2.25);
  \node[font=\small, anchor=east] at (axis cs:16.3,2.4) {$\ell+\varepsilon$};
  \node[font=\small, anchor=east] at (axis cs:16.3,1.6) {$\ell-\varepsilon$};
\end{axis}
\end{tikzpicture}

{\small Convergence of $u_n = 2 + \frac{(-1)^n}{n}$ to $\ell = 2$: given
$\varepsilon > 0$, all terms from the index $N$ on lie in the band
$\intcc{\ell-\varepsilon}{\ell+\varepsilon}$.}
\end{omfigure}

\begin{definition}[Divergence to infinity]\label{def:g12:seq:infinity}
The \omterm{def:g12:seq:sequence}{sequence} $(u_n)$ \emph{tends to $+\infty$} if for every $A \in \R$, there
exists $N \in \N$ such that $u_n \geq A$ for all $n \geq N$. We write
$\lim\limits_{n\to+\infty} u_n = +\infty$; the definition of
$\lim u_n = -\infty$ is analogous. A \omterm{def:g12:seq:sequence}{sequence} that does not converge is said
to \emph{diverge}.
\end{definition}

\begin{remark}
A \omterm{def:g12:seq:sequence}{sequence} can diverge without tending to $\pm\infty$: the \omterm{def:g12:seq:sequence}{sequence}
$u_n = (-1)^n$ takes only the values $1$ and $-1$ and has no limit.
\end{remark}

\begin{proposition}[Uniqueness of the limit]\label{prop:g12:seq:uniqueness}
If $(u_n)$ \omterm{def:g12:seq:limit}{converges}, its limit is unique.
\end{proposition}

\begin{proof}
Suppose $u_n \to \ell$ and $u_n \to \ell'$ with $\ell \neq \ell'$, say
$\ell < \ell'$. Set $\varepsilon = \frac{\ell' - \ell}{3} > 0$. From some
index on, $\abs{u_n - \ell} \leq \varepsilon$ and
$\abs{u_n - \ell'} \leq \varepsilon$, hence
\[
\ell' - \ell \leq \abs{\ell' - u_n} + \abs{u_n - \ell}
\leq 2\varepsilon = \tfrac{2}{3}(\ell' - \ell) < \ell' - \ell,
\]
a contradiction.
\end{proof}

\begin{proposition}[Operations on limits]\label{prop:g12:seq:operations}
Let $(u_n)$ and $(v_n)$ be \omterm{def:g12:seq:sequence}{sequences} with limits $\ell$ and $\ell'$ (finite or
infinite). Then, whenever the right-hand side is not an indeterminate form,
\[
\lim (u_n + v_n) = \ell + \ell', \qquad
\lim (u_n v_n) = \ell\,\ell', \qquad
\lim \frac{u_n}{v_n} = \frac{\ell}{\ell'}.
\]
The indeterminate forms are
$(+\infty) + (-\infty)$, $0 \times \infty$, $\frac{\infty}{\infty}$ and
$\frac{0}{0}$.
\end{proposition}

\begin{proof}
We prove the sum rule for finite limits; the other cases are similar and left
as exercises. Let $\varepsilon > 0$. There exist $N_1, N_2$ such that
$\abs{u_n - \ell} \leq \varepsilon/2$ for $n \geq N_1$ and
$\abs{v_n - \ell'} \leq \varepsilon/2$ for $n \geq N_2$. For
$n \geq \max(N_1, N_2)$, the triangle inequality gives
\[
\abs{(u_n + v_n) - (\ell + \ell')}
\leq \abs{u_n - \ell} + \abs{v_n - \ell'} \leq \varepsilon. \qedhere
\]
\end{proof}

\begin{method}[Lifting an indeterminate form]\label{met:g12:seq:indeterminate}
Faced with an indeterminate form, factor out the dominant term. For instance
\[
n^2 - n = n^2\left(1 - \tfrac{1}{n}\right) \xrightarrow[n\to+\infty]{} +\infty,
\qquad
\frac{2n^2+1}{n^2 - n} = \frac{2 + 1/n^2}{1 - 1/n}
\xrightarrow[n\to+\infty]{} 2 .
\]
\end{method}

\section{Convergence theorems}

\begin{theorem}[Comparison and squeeze theorems]\label{thm:g12:seq:squeeze}
Let $(u_n)$, $(v_n)$, $(w_n)$ be \omterm{def:g12:seq:sequence}{sequences}.
\begin{enumerate}
  \item If $u_n \leq v_n$ from some index on and $u_n \to +\infty$, then
        $v_n \to +\infty$.
  \item \emph{(Squeeze theorem)}\index{squeeze theorem} If
        $u_n \leq v_n \leq w_n$ from some index on and $(u_n)$ and $(w_n)$
        both converge to the same limit $\ell$, then $(v_n)$ \omterm{def:g12:seq:limit}{converges}
        to~$\ell$.
\end{enumerate}
\end{theorem}

\begin{proof}
\emph{1.} Let $A \in \R$. Since $u_n \to +\infty$, there is $N$ with
$u_n \geq A$ for $n \geq N$; enlarging $N$ if needed, $v_n \geq u_n \geq A$
for $n \geq N$.

\emph{2.} Let $\varepsilon > 0$. From some index on, both
$\ell - \varepsilon \leq u_n$ and $w_n \leq \ell + \varepsilon$ hold, hence
$\ell - \varepsilon \leq u_n \leq v_n \leq w_n \leq \ell + \varepsilon$, that
is $\abs{v_n - \ell} \leq \varepsilon$.
\end{proof}

\begin{example}
For all $n \geq 1$, $-\frac{1}{n} \leq \frac{(-1)^n}{n} \leq \frac{1}{n}$, and
both bounds tend to $0$; hence $\frac{(-1)^n}{n} \to 0$.
\end{example}

\begin{theorem}[Monotone convergence theorem]\label{thm:g12:seq:monotone}
An \omterm{def:g12:seq:monotonic}{increasing} \omterm{def:g12:seq:sequence}{sequence} that is \omterm{def:g12:seq:bounded}{bounded above} \omterm{def:g12:seq:limit}{converges}. A \omterm{def:g10:functions:variations}{decreasing} \omterm{def:g12:seq:sequence}{sequence}
that is \omterm{def:g12:seq:bounded}{bounded below} \omterm{def:g12:seq:limit}{converges}. An \omterm{def:g12:seq:monotonic}{increasing} \omterm{def:g12:seq:sequence}{sequence} that is not \omterm{def:g12:seq:bounded}{bounded
above} tends to $+\infty$.
\end{theorem}

\begin{proof}[Partial proof]
We prove the third statement. Let $(u_n)$ be \omterm{def:g12:seq:monotonic}{increasing} and not \omterm{def:g12:seq:bounded}{bounded
above}, and let $A \in \R$. Since $A$ is not an upper bound, there exists $N$
with $u_N \geq A$; by monotonicity, $u_n \geq u_N \geq A$ for all $n \geq N$.
Hence $u_n \to +\infty$.

The two convergence statements rely on the least upper bound property of
$\R$; they are \emph{admitted at this level} (and proved in the first year of
university).
\end{proof}

\begin{remark}
The theorem guarantees the \emph{existence} of the limit but does not give
its value. An \omterm{def:g12:seq:monotonic}{increasing} \omterm{def:g12:seq:sequence}{sequence} \omterm{def:g12:seq:bounded}{bounded above} by $M$ \omterm{def:g12:seq:limit}{converges} to some
$\ell \leq M$, not necessarily to $M$.
\end{remark}

\begin{theorem}[Limit of geometric sequences]\label{thm:g12:seq:geometric}
Let $q \in \R$.
\begin{enumerate}
  \item If $q > 1$, then $q^n \to +\infty$.
  \item If $q = 1$, then $q^n \to 1$.
  \item If $\abs{q} < 1$, then $q^n \to 0$.
  \item If $q \leq -1$, then $(q^n)$ diverges and has no limit.
\end{enumerate}
\end{theorem}

\begin{proof}
\emph{1.} Write $q = 1 + a$ with $a > 0$. Bernoulli's inequality
(\cref{ex:g12:seq:bernoulli}) gives $q^n \geq 1 + na \to +\infty$, and we
conclude by comparison (\cref{thm:g12:seq:squeeze}).

\emph{2.} Immediate.

\emph{3.} If $q = 0$ the claim is clear. Otherwise $\abs{q} < 1$ gives
$1/\abs{q} > 1$, so $(1/\abs{q})^n \to +\infty$ by point 1, hence
$\abs{q}^n \to 0$, and $-\abs{q}^n \leq q^n \leq \abs{q}^n$ allows us to
conclude by the squeeze theorem.

\emph{4.} For $q \leq -1$, $(q^{2n})$ takes values $\geq 1$ while
$(q^{2n+1})$ takes values $\leq -1$: no single limit can attract both
subsequences.
\end{proof}

\begin{omfigure}
\begin{tikzpicture}
\begin{axis}[omaxis,
  width=10cm, height=5.6cm,
  xlabel={$n$}, ylabel={$q^n$},
  xmin=-0.5, xmax=13, ymin=-1.3, ymax=4.3,
  xtick={5,10}, ytick={-1,1}]
  \addplot[omThm, only marks, mark=*, mark size=1.5pt,
    samples at={0,...,6}] {1.25^x};
  \addplot[omDef, only marks, mark=*, mark size=1.5pt,
    samples at={0,...,12}] {0.75^x};
  \addplot[omProp, only marks, mark=*, mark size=1.5pt,
    samples at={0,...,12}] {(-0.8)^x};
  \node[omThm, font=\small, anchor=west] at (axis cs:6.3,3.7) {$q=1.25$};
  \node[omDef, font=\small, anchor=south west] at (axis cs:3.4,0.55) {$q=0.75$};
  \node[omProp, font=\small, anchor=north] at (axis cs:7,-0.75) {$q=-0.8$};
\end{axis}
\end{tikzpicture}

{\small The three behaviours of $(q^n)$: divergence to $+\infty$ for
$q > 1$ (red), convergence to $0$ for $\abs q < 1$ (blue), and damped
oscillation --- still convergence to $0$ --- for $-1 < q < 0$ (orange).}
\end{omfigure}

\begin{method}[Recurrent sequences $u_{n+1} = f(u_n)$]\label{met:g12:seq:recurrent}
To study a \omterm{def:g12:seq:sequence}{sequence} defined by $u_{n+1} = f(u_n)$:
\begin{enumerate}
  \item prove by induction that $(u_n)$ stays in an \omterm{def:g10:numbers:interval}{interval} $I$ on which $f$
        is well behaved (and, often, that $(u_n)$ is \omterm{def:g12:seq:monotonic}{monotonic});
  \item deduce convergence from the monotone convergence theorem;
  \item pass to the limit in the relation $u_{n+1} = f(u_n)$: if $f$ is
        continuous and $u_n \to \ell \in I$, then $\ell$ satisfies
        $f(\ell) = \ell$ (see \cref{ch:g12:limcont}); solve this \omterm{def:g10:algebra:equation}{equation} and
        select the right \omterm{def:g11:quad:discriminant}{root}.
\end{enumerate}
\end{method}

\begin{omfigure}
\begin{tikzpicture}
\begin{axis}[omaxis,
  width=8.4cm, height=6.4cm,
  xmin=-0.4, xmax=3.1, ymin=-0.4, ymax=2.7,
  xtick={0,1.414,1.848}, xticklabels={$u_0$,$u_1$,$u_2$}, ytick=\empty]
  \addplot[omProp, dashed, domain=-0.3:2.6] {x};
  \addplot[omDef, thick, domain=-0.4:2.9] {sqrt(x+2)};
  \draw[omThm, thick] (axis cs:0,0) -- (axis cs:0,1.414)
    -- (axis cs:1.414,1.414) -- (axis cs:1.414,1.848)
    -- (axis cs:1.848,1.848) -- (axis cs:1.848,1.961)
    -- (axis cs:1.961,1.961);
  \draw[black, dashed, thin] (axis cs:1.414,0) -- (axis cs:1.414,1.414);
  \draw[black, dashed, thin] (axis cs:1.848,0) -- (axis cs:1.848,1.848);
  \fill (axis cs:2,2) circle (1.5pt)
    node[below right, font=\small] {$\ell = 2$};
  \node[omDef, font=\small, above left] at (axis cs:2.75,2.18)
    {$y = f(x)$};
  \node[omProp, font=\small, anchor=north west] at (axis cs:1.08,0.82)
    {$y = x$};
\end{axis}
\end{tikzpicture}

{\small Staircase construction for $u_{n+1} = \sqrt{u_n + 2}$, $u_0 = 0$
(\cref{exo:g12:seq:6}): each vertical step reads $f(u_n)$ on the curve,
each horizontal step carries it back through $y = x$. The \omterm{def:g12:seq:sequence}{sequence} climbs
to the \omterm{pb:g10:functions:1}{fixed point} $\ell = 2$, where the curve meets the line.}
\end{omfigure}

\begin{example}\label{ex:g12:seq:sqrt2}
Let $u_0 = 2$ and $u_{n+1} = \frac{1}{2}\left(u_n + \frac{2}{u_n}\right)$.
One checks by induction that $u_n \geq \sqrt{2}$ for all $n$ (the inequality
$\frac{1}{2}(x + 2/x) \geq \sqrt{2}$ for $x>0$ is equivalent to
$(x - \sqrt2)^2 \geq 0$), then that $(u_n)$ is \omterm{def:g10:functions:variations}{decreasing}, since
\[
u_{n+1}-u_n=\frac{2-u_n^2}{2u_n}\leq 0 .
\]
\omterm{def:g10:functions:variations}{Decreasing} and \omterm{def:g12:seq:bounded}{bounded below}, $(u_n)$
\omterm{def:g12:seq:limit}{converges} to some $\ell \geq \sqrt{2}$, which must satisfy
$\ell = \frac{1}{2}(\ell + 2/\ell)$, i.e.\ $\ell^2 = 2$. Hence
$u_n \to \sqrt{2}$. This is Heron's algorithm, already used by the
Babylonians; its convergence is extremely fast ($u_3$ already gives $\sqrt 2$
to eight decimal places).
\end{example}

\section{Exercises}

\begin{exercise}[$\star$]\label{exo:g12:seq:1}
Prove by induction that for all $n \in \N$,
\[
1^2 + 2^2 + \dots + n^2 = \frac{n(n+1)(2n+1)}{6}.
\]
\end{exercise}

\begin{exercise}[$\star$]\label{exo:g12:seq:2}
Study the monotonicity of the \omterm{def:g12:seq:sequence}{sequences} defined for $n \geq 1$ by
\[
a_n = \frac{n+1}{n}, \qquad
b_n = \frac{2^n}{n}, \qquad
c_n = n^2 - 10n .
\]
\end{exercise}

\begin{exercise}[$\star$]\label{exo:g12:seq:3}
Compute the limits of the \omterm{def:g12:seq:sequence}{sequences} with general terms
\[
u_n = \frac{3n^2 - n + 1}{2n^2 + 5}, \qquad
v_n = \sqrt{n+1} - \sqrt{n}, \qquad
w_n = \frac{2^n - 3^n}{3^n + 1}.
\]
\end{exercise}

\begin{exercise}[$\star$]\label{exo:g12:seq:4}
Let $(u_n)$ be the \omterm{def:g11:seq:arithmetic}{arithmetic sequence} with $u_0 = 5$ and \omterm{def:g11:seq:arithmetic}{common difference}
$r = 3$, and $(v_n)$ the \omterm{def:g11:seq:geometric}{geometric sequence} with $v_0 = 8$ and \omterm{def:g11:seq:geometric}{common ratio}
$q = \frac{1}{2}$. Compute $u_n$, $v_n$, $\sum_{k=0}^{n} u_k$ and
$\sum_{k=0}^{n} v_k$, and the limits of all four expressions as
$n \to +\infty$.
\end{exercise}

\begin{exercise}[$\star\star$]\label{exo:g12:seq:5}
Using the squeeze theorem, compute
\[
\lim_{n\to+\infty} \frac{n + \cos n}{n + 1}
\qquad\text{and}\qquad
\lim_{n\to+\infty} \frac{n!}{n^n},
\]
where $n! = 1 \times 2 \times \dots \times n$. For the second limit, bound
$\frac{n!}{n^n}$ by a term of a \omterm{def:g11:seq:geometric}{geometric sequence}.
\end{exercise}

\begin{exercise}[$\star\star$]\label{exo:g12:seq:6}
Let $u_0 = 0$ and $u_{n+1} = \sqrt{u_n + 2}$ for all $n \in \N$.
\begin{enumerate}
  \item Prove by induction that $0 \leq u_n \leq 2$ for all $n$.
  \item Show that $(u_n)$ is \omterm{def:g12:seq:monotonic}{increasing}.
  \item Deduce that $(u_n)$ \omterm{def:g12:seq:limit}{converges} and determine its limit.
\end{enumerate}
\end{exercise}

\begin{exercise}[$\star\star$]\label{exo:g12:seq:7}
A patient takes a dose of $1$ unit of a drug every morning. During each
24-hour period, the body eliminates $40\%$ of the drug present. Let $u_n$ be
the quantity of drug in the body just after the dose on day $n$, so that
$u_0 = 1$.
\begin{enumerate}
  \item Justify that $u_{n+1} = 0.6\,u_n + 1$.
  \item Let $v_n = u_n - 2.5$. Show that $(v_n)$ is \omterm{def:g12:seq:arith-geom}{geometric} and deduce an
        explicit formula for $u_n$.
  \item Determine the long-term quantity of drug in the body.
\end{enumerate}
\end{exercise}

\begin{exercise}[$\star\star$]\label{exo:g12:seq:8}
Let $(u_n)$ be defined by $u_0 = 3$ and $u_{n+1} = \frac{4u_n - 1}{u_n + 2}$.
\begin{enumerate}
  \item Show by induction that $u_n > 1$ for all $n \in \N$.
  \item Show that $v_n = \dfrac{1}{u_n - 1}$ defines an \omterm{def:g11:seq:arithmetic}{arithmetic sequence}.
  \item Deduce explicit formulas for $v_n$ and $u_n$, and the limit of
        $(u_n)$.
\end{enumerate}
\end{exercise}

\begin{exercise}[$\star\star\star$]\label{exo:g12:seq:9}
For $n \geq 1$, let
$H_n = 1 + \frac{1}{2} + \frac{1}{3} + \dots + \frac{1}{n}$.
\begin{enumerate}
  \item Show that for all $n \geq 1$, $H_{2n} - H_n \geq \frac{1}{2}$.
  \item Deduce that $H_{2^k} \geq 1 + \frac{k}{2}$ for all $k \in \N$, and
        conclude that $H_n \to +\infty$.
\end{enumerate}
\end{exercise}

\begin{exercise}[$\star\star\star$]\label{exo:g12:seq:10}
\emph{(Adjacent \omterm{def:g12:seq:sequence}{sequences}.)} Two \omterm{def:g12:seq:sequence}{sequences} $(a_n)$ and $(b_n)$ are
\emph{adjacent} if $(a_n)$ is \omterm{def:g12:seq:monotonic}{increasing}, $(b_n)$ is \omterm{def:g10:functions:variations}{decreasing}, and
$b_n - a_n \to 0$.
\begin{enumerate}
  \item Show that for all $n$, $a_n \leq b_n$. (Hint: study the monotonicity
        of $(b_n - a_n)$.)
  \item Show that adjacent \omterm{def:g12:seq:sequence}{sequences} both converge, to the same limit.
  \item Application: show that the \omterm{def:g12:seq:sequence}{sequences}
        $a_n = \sum_{k=0}^{n} \frac{1}{k!}$ and
        $b_n = a_n + \frac{1}{n \cdot n!}$ ($n \geq 1$) are adjacent. (Their
        common limit is the number $\eu$, studied in \cref{ch:g12:exp}.)
\end{enumerate}
\end{exercise}

\section{Problem: Heron's sequence, judged at last}

\begin{problem}[{Weekend problem --- induction certifies, monotone
convergence sentences, and the two-thousand-year-old recipe for
$\sqrt2$ finally gets its proof (with Gauss's marvellous mean for
dessert)}]
\label{pb:g12:seq:1}
Three times this series has met Heron's recipe --- average the
guess with $2/\text{guess}$ --- and three times it could only
\emph{observe} that the recipe works. This chapter at last owns
the instruments of judgment: induction
(\cref{thm:g12:seq:induction}), the monotone convergence theorem
(\cref{thm:g12:seq:monotone}), and limits of recurrences. The
verdict, and the certified speed, occupy the heart of this
problem; around it, induction's classic traps, the slowest
divergence in mathematics, and the fastest convergence Gauss ever
found.

\medskip
\noindent\textbf{Part I --- Induction warm-ups.}
\begin{enumerate}
  \item Prove by induction:
        $1 + 3 + 5 + \dots + (2n - 1) = n^2$ (the staircase of
        odd numbers, drawn in the Middle School volume, now
        certified).
  \item Prove by induction that $2^n > n$ for every
        $n \in \N$.
  \item Prove Bernoulli's inequality by induction: for
        $x \geq 0$ and $n \in \N$,
        $(1 + x)^n \geq 1 + nx$.
  \item The classic trap: ``all marbles have the same color ---
        true for one marble; and if any $n$ marbles are always
        monochrome, then among $n + 1$ marbles the first $n$
        share a color, the last $n$ share a color, so all
        $n + 1$ do.'' Every child knows the conclusion is
        absurd: find the exact step where the induction breaks.
  \item Prove by induction that $4^n - 1$ is divisible by $3$
        for every $n \in \N$.
\end{enumerate}

\medskip
\noindent\textbf{Part II --- The trial of Heron.}
Let $x_0 = 2$ and $x_{n+1} = \dfrac12\left(x_n +
\dfrac{2}{x_n}\right)$.
\begin{enumerate}[resume]
  \item Compute $x_1$, $x_2$, $x_3$ as exact fractions (old
        friends).
  \item Prove the key identity
        \[
        x_{n+1}^2 - 2 = \left(\frac{x_n^2 - 2}{2x_n}\right)^{\!2}
        \geq 0,
        \]
        and deduce by induction that $x_n > 0$ and
        $x_n^2 > 2$ for every $n$.
  \item Show that $(x_n)$ is strictly \omterm{def:g10:functions:variations}{decreasing} (compute
        $x_{n+1} - x_n$ and use question 7).
  \item Invoke the monotone convergence theorem: why does
        $(x_n)$ converge to some limit $L \geq 1$?
  \item Identify the limit: pass the recurrence to the limit
        (\cref{prop:g12:seq:operations}) and conclude
        $L = \sqrt2$. State the historical verdict: after two
        thousand years of faithful service, Heron's recipe is
        \emph{proved} to converge.
  \item The certified speed: with $e_n = x_n - \sqrt2$, prove
        \[
        e_{n+1} = \frac{e_n^2}{2 x_n} ,
        \]
        and deduce $e_{n+1} \leq \frac{e_n^2}{2\sqrt2}$: the
        error is squared at each step --- the digit-doubling
        observed since the Middle School volume, now a theorem.
  \item Confirm numerically: compute $e_0, e_1, e_2, e_3$ (from
        question 6) and check that each $\frac{e_{n+1}}{e_n^2}$
        is close to $\frac{1}{2x_n}$.
\end{enumerate}

\medskip
\noindent\textbf{Part III --- The slowest divergence.}
\begin{enumerate}[resume]
  \item \Cref{exo:g12:seq:9} proved
        $H_{2^k} \geq 1 + \frac k2$ for the harmonic sums. How
        many terms guarantee $H_n > 10$? (A power of two will
        do; marvel at its size.)
  \item By contrast, the \omterm{def:g12:seq:arith-geom}{geometric} sums
        $1 + \frac12 + \frac14 + \dots + \frac{1}{2^n} =
        2 - \frac{1}{2^n}$ converge to $2$
        (\cref{thm:g12:seq:geometric}): the chocolate-bar
        intuition of the Middle School volume, finally a limit
        statement. Write out the two-line proof.
  \item Between the two: show that the sums
        $S_n = 1 + \frac{1}{4} + \frac{1}{9} + \dots +
        \frac{1}{n^2}$ converge, by bounding
        $\frac{1}{k^2} \leq \frac{1}{k(k-1)} =
        \frac{1}{k-1} - \frac{1}{k}$ (for $k \geq 2$),
        telescoping, and applying monotone convergence. (The
        limit, $\frac{\pi^2}{6}$, is one of Euler's miracles,
        proved in the university volumes.)
  \item State the moral of questions 13--15 in two sentences:
        what does ``the terms tend to $0$'' decide about the
        convergence of the sums --- and what does it not?
\end{enumerate}

\medskip
\noindent\textbf{Part IV --- Gauss's arithmetic--geometric
\omterm{def:g11:stat:mean}{mean}.}
Let $a_0 = 1$, $b_0 = 2$, and
\[
a_{n+1} = \sqrt{a_n b_n},
\qquad
b_{n+1} = \frac{a_n + b_n}{2} .
\]
\begin{enumerate}[resume]
  \item Compute $a_1, b_1, a_2, b_2$ (five decimals). What do
        you observe about the speed?
  \item Show that $a_n \leq b_n$ for every $n$ (the
        arithmetic--geometric inequality, met throughout this
        series), that $(a_n)$ increases and $(b_n)$ decreases.
  \item Show that
        $b_{n+1} - a_{n+1} \leq \frac{b_n - a_n}{2}$
        (factor $b_{n+1} - a_{n+1} =
        \frac{(\sqrt{b_n} - \sqrt{a_n})^2}{2}$ and compare),
        and conclude with \cref{exo:g12:seq:10} that the two
        \omterm{def:g12:seq:sequence}{sequences} are adjacent: they share a common limit
        $M(1, 2)$, the \emph{arithmetic--geometric \omterm{def:g11:stat:mean}{mean}}.
  \item Compute $M(1, 2)$ to six decimals (how many iterations
        did you need?). On 30 May 1799, Gauss computed
        $M(1, \sqrt2)$ to eleven decimals, recognized
        $\frac{\pi}{M(1,\sqrt2)}$ as a known integral, and
        wrote that a ``new field of analysis'' had opened ---
        it had: elliptic integrals, told in the university
        volumes. Close with the observed convergence speeds of
        this problem, slowest to fastest.
\end{enumerate}
\end{problem}
